\chapter{Rare Event Simulation}
\label{ch:rareevent}

Counting queens is estimating $\PP_p(S = 0)$ where $p$ is uniform on permutations, and $\PP_p(S = 0) \approx e^{-0.942 N}$. This is a rare event probability, and there is a mature theory of estimating such probabilities, grown up in reliability, queueing, finance and particle physics. This chapter reviews it from the vertical point of view: every method in the literature is a way of reweighting or rearranging the quantile function\index{quantile function} so that the rare region is reached, and the methods differ in what they assume about the structure of that region.

\section{The relative error criterion}

An estimator $\hat\rho_N$ of a rare event probability $\rho_N \to 0$ has bounded relative error if $\Var(\hat\rho_N) / \rho_N^2$ stays bounded as $N \to \infty$, and is logarithmically efficient if $\log \EE[\hat\rho_N^2] / \log \rho_N^2 \to 1$. The naive estimator has relative variance $1/\rho_N$ per sample, exponential in $N$. The goal of every rare event method is a polynomial number of samples for bounded relative error, and in the complexity vocabulary this is exactly an FPRAS\index{FPRAS}: a $(1 \pm \epsilon)$ estimate with constant probability from polynomially many samples of polynomial cost each.

\section{Importance sampling and exponential tilting}

Importance sampling draws from a density $g$ instead of $p$ and weights by $p/g$:
\[
\hat\rho = \frac{1}{M} \sum_{i=1}^{M} \Ind_{\{S(x^{(i)}) = 0\}} \frac{p(x^{(i)})}{g(x^{(i)})} , \qquad x^{(i)} \sim g .
\]
The estimator is unbiased for any $g$ positive on the event and has zero variance when $g = g^* = p(\cdot \mid S = 0)$, the conditional law on the event. The optimal density is the thing we are trying to count, which is no help directly, and the practical question is how close one can get to it with a tractable $g$.

Exponential tilting takes $g_\gamma(x) \propto p(x) e^{-S(x)/\gamma}$, the Boltzmann distribution at temperature $\gamma$ (inverse temperature $1/\gamma$). Siegmund showed in 1976 that for sums of independent variables the tilt that moves the mean of $S$ to the boundary of the rare event is asymptotically optimal, and the theory of large deviations, through Cram\'er's theorem and its generalisations by Varadhan and by Bucklew, explains why: the rare event is dominated by the most likely way of achieving it, and the tilt that makes that way typical is the one to use. For the queens the tilted law $g_\gamma$ is the Boltzmann distribution of Chapter~\ref{ch:rare} and the tilt $\gamma \to 0$ moves the mean energy to zero. Simulated annealing\index{simulated annealing} used as an estimator is exponential tilting with the tilt varied along a schedule, and the product estimator is the resulting ratio of normalising constants.

In the vertical coordinate, tilting changes the quantile function. Under $g_\gamma$ the step of the staircase at energy $s$ has width proportional to $N(s) e^{-s/\gamma}$ rather than $N(s)$, so decreasing $\gamma$ moves mass towards the left, where the solutions are, and the walk from $s = 1$ to $s = 0$ becomes shorter. The limit $\gamma \to 0$ puts all the mass on the ground step and the walk has nothing to do, but the normalising constant of $g_\gamma$ is then exactly the unknown count. Tilting trades walk length for normalisation; the ladder of tilts is the ladder of levels.

\section{Splitting}

Kahn and Harris in 1951 proposed, for neutron transport, that trajectories which approach a rare region be split into several copies, each weighted by the reciprocal of the number of copies, so that the rare region is reached by many trajectories while the estimator stays unbiased. Garvels, Kroese and others developed the method into multilevel splitting: choose levels $m_T > m_{T-1} > \dots > m_0$, run particles until they reach the next level down, split those that do, and estimate the probability of the final level as the product of the level crossing probabilities. This is precisely the product estimator\index{product estimator} of Chapter~\ref{ch:product}, and the choice of levels at fixed quantiles of the current population, as in the adaptive multilevel splitting of C\'erou and Guyader, is precisely the quantile ladder of split sampling.

The analysis of multilevel splitting shows that the relative variance is minimised when the level crossing probabilities are equal, and that the optimal crossing probability is a constant near $0.2$ (Lagnoux~\cite{Lagnoux}, 2006), so that the number of levels is of order $\log(1/\rho)$. This is the well balanced schedule of Chapter~\ref{ch:product} derived from the splitting side. Lagnoux and Glasserman, Heidelberger, Shahabuddin and Zajic analysed the method under the assumption that particles at a level are distributed as the conditional law, and showed that under that assumption the sample complexity is $O(\log(1/\rho)^2 / \epsilon^2)$. The assumption is the mixing assumption, and Chapter~\ref{ch:mixing} is about when it holds.

\section{The zero variance estimator and its approximations}

Every unbiased estimator of $\rho = \EE_p[\Ind_{\{S = 0\}}]$ has an associated importance density, and the zero variance density is $p(\cdot \mid S = 0)$. One can approach it by a sequence of densities $g_t = p(\cdot \mid S \le m_t)$, which is splitting; by a parametric family optimised in Kullback Leibler divergence, which is the cross entropy\index{cross entropy method} method; or by a product of marginals, which is the mean field or variational approximation. Each is a different way of approximating the Lorenz curve of $\Ind_{\{S=0\}}$ under $p$, and the quality of the approximation is the relative variance.

The cross entropy method deserves a second look in the vertical coordinate. With $g_\eta$ an exponential family and the elite sample the fraction $\kappa$ of draws with smallest $S$, the update sets $\eta$ to match the sufficient statistics of the elite sample. For the queens with $g_\eta$ a product of row marginals this makes $\eta_{ij}$ large where queens in the elite tend to sit. The one point marginals of the uniform law on solutions are the finite $N$ precursors of the queenon\index{queenon} of Simkin's theorem, but we do not claim that the method learns them: in combinatorial optimisation the cross entropy iteration is known to converge to a degenerate law concentrated on a single good configuration \cite{RubinsteinKroeseCE}. Whether the product family is rich enough for the relative variance to stay bounded is doubtful, since solutions are far from product measures, and the run of Chapter~\ref{ch:ce} finds no solutions for $9 \le N \le 12$.

\section{Diaconis and Holmes}

The paper of Diaconis\index{Diaconis and Holmes} and Holmes in 1995 is the direct ancestor of our work. It treats three examples at the interface of statistics, computer science and statistical mechanics, and the queens is one. Their approach to counting solutions is to consider the Gibbs distribution with energy the number of attacking pairs, to run a Metropolis\index{Metropolis algorithm} chain, and to estimate the number of solutions from the proportion of time the chain spends at zero energy, with the normalising constant obtained by a telescoping product over temperatures. They observe that the problem of counting is thereby replaced by the problem of estimating a rare event probability, which is the first key insight of our 2024 paper, and they note that the Metropolis chain mixes slowly at low temperature. Their remarks on the relation between the sampling and counting problems anticipate the Jerrum Valiant\index{Valiant, Leslie} Vazirani theorem as applied to the queens and our corollary in Chapter~\ref{ch:prp} is the formal version of what they said informally.

\section{Large deviations and the queenon}

Simkin\index{Simkin's theorem}'s theorem is a large deviations result. The probability that a uniform point of the $N^N$ space has zero attacking pairs is $e^{-\alpha N + o(N)}$, and for a uniform permutation it is $e^{-(\alpha - 1) N + o(N)}$, and the rate constant is characterised by a variational problem over densities on the unit square: the queenon $f$ maximises a functional that balances an entropy term against the constraint that no two queens attack. The analogy with Sanov's theorem is close. Sanov says that the probability that the empirical measure of a sample is near a given measure $\nu$ decays like $e^{-N D(\nu \| \mu)}$; Simkin says that the probability that a random permutation is a queens solution decays at a rate given by the entropy of the most entropic queenon.

The queenon is the large deviations tilt: the conditional law of a uniform permutation given that it is a solution converges, in the sense of limit objects, to the queenon, and the optimal importance density is $p(\cdot \mid S = 0)$, whose one point marginals are the finite $N$ precursors of the queenon. An importance density built from the queenon alone would not be $p(\cdot \mid S = 0)$, since the queenon records only one point marginals, and we do not claim that it would have bounded relative error, which requires $g \ge c \, p(\cdot \mid S = 0)$ on the event uniformly in $N$. The optimal density itself cannot be normalised, because the normalisation is the count, and this is the usual circularity of rare event simulation: the optimal density is the answer. The vertical walk breaks the circle by tilting in stages and normalising each stage against the previous one.

\section{Rare events in the quantile coordinate}

We can now state the general principle. Let $Y = L(X)$ be the likelihood ordinate, with quantile function $\Lambda$, and suppose the quantity of interest is $Z = \int_0^1 \Lambda(s) \, ds$ with $\Lambda$ concentrated near $s = 0$, so that $Z$ is dominated by a rare region of the quantile axis. Importance sampling in the quantile coordinate draws $s \sim \Omega$ and weights by $\Lambda(s) / \Omega(s)$; the optimal $\Omega$ is $\Lambda / Z$ and the rare event methods are different schemes for approximating it. Tilting multiplies $\Lambda$ by a function of the level and renormalises; splitting discretises $\Omega$ onto a ladder and estimates each rung by conditioned draws; cross entropy approximates $\Omega$ within a parametric family. All are the same thing seen on the unit interval, and the vertical likelihood\index{vertical likelihood} framework is simply the recognition that the unit interval is where the problem lives.
